\fancyhead[LO]{\makebox[\textwidth][r]{Chapter 2. Background and Notation}}
\chapter{Background and Notation} \label{chapter:background}

This chapter establishes the notation and mathematical background used throughout the dissertation and briefly positions prior work.
\sref{sec:ch2_notation} introduces vectors, matrices, frames, and Lie groups used in later chapters.
\sref{sec:ch2_estimation} reviews probabilistic state estimation, factor graph optimization, and the role of measurement covariance.
\sref{sec:ch2_slam} formalizes the SLAM and visual-inertial estimation problem.
\sref{sec:ch2_vo} and \sref{sec:ch2_imu_preint} review stereo visual odometry and IMU integration with uncertainty propagation.
\sref{sec:ch2_related_works} provides a concise related-work overview tied to the three parts of the thesis.

\section{Notation}
\label{sec:ch2_notation}

\paragraph{Vectors, matrices, and sets.}
Bold lowercase letters ($\mathbf{x}$) denote vectors, bold uppercase letters ($\mathbf{R}$) denote matrices, and calligraphic letters ($\mathcal{X}$) denote sets.
$\mathbf{I}_n$ is the $n \times n$ identity matrix.
$\mathbf{0}$ denotes the zero vector or matrix; the dimension is clear from context.
Time indices appear as subscripts ($\mathbf{x}_k$ at time $k$) and ranges as $\mathbf{x}_{1:T} = \{\mathbf{x}_1,\dots,\mathbf{x}_T\}$.

\paragraph{Reference frames and rigid-body transformations.}
We use $\mathbf{R} \in \mathrm{SO}(3)$ for rotation matrices and $\mathbf{T} \in \mathrm{SE}(3)$ for rigid-body transformations.
Superscripts and subscripts indicate reference frames; $\mathbf{T}^{a}_{b}$ transforms vectors from frame $\{b\}$ to frame $\{a\}$, and $\mathbf{p}^a$ denotes a 3D point expressed in frame $\{a\}$.
Common frames in this thesis are the world frame $\{w\}$, the body or IMU frame $\{b\}$, and the camera frame $\{c\}$.

\paragraph{Lie groups and operators.}
The operator $[\,\cdot\,]_\times$ constructs the skew-symmetric matrix from a 3-vector.
$\mathrm{Exp}(\cdot)$ and $\mathrm{Log}(\cdot)$ denote the exponential and logarithmic maps between Lie groups and their algebras, with $\mathrm{Exp}: \mathbb{R}^3 \to \mathrm{SO}(3)$ and the corresponding $\mathrm{SE}(3)$ maps used for rigid-body increments.
The boxplus operator $\mathbf{R} \boxplus \delta\bm{\phi} = \mathbf{R}\,\mathrm{Exp}(\delta\bm{\phi})$ defines local perturbations on $\mathrm{SO}(3)$.

\paragraph{Probability and uncertainty.}
$\mathcal{N}(\bm{\mu},\bm{\Sigma})$ denotes a multivariate Gaussian with mean $\bm{\mu}$ and covariance $\bm{\Sigma}$.
The squared Mahalanobis norm of a residual $\mathbf{r}$ under covariance $\bm{\Sigma}$ is
\begin{equation}
\label{eq:ch2:mahalanobis}
\|\mathbf{r}\|^2_{\bm{\Sigma}} = \mathbf{r}^{\top} \bm{\Sigma}^{-1} \mathbf{r}.
\end{equation}
We write $\bm{\Sigma}_k$ for the covariance of the $k$-th measurement and $\bm{\Sigma}_k(\mathcal{Z};\theta)$ when the covariance is data-conditioned through learned parameters $\theta$.

\section{Probabilistic State Estimation and Optimization}
\label{sec:ch2_estimation}

\subsection{MAP Estimation}
\label{sec:ch2_map}

State estimation in this thesis is formulated as maximum a posteriori (MAP) inference.
Given a state $\mathcal{X}$ and a set of measurements $\mathcal{Z} = \{\mathbf{z}_1,\dots,\mathbf{z}_m\}$ with conditional likelihoods $p(\mathbf{z}_k \mid \mathcal{X}_k) \propto \exp\!\left(-\tfrac{1}{2}\|\mathbf{r}_k(\mathcal{X}_k,\mathbf{z}_k)\|^2_{\bm{\Sigma}_k}\right)$ and prior $p(\mathcal{X})$, the MAP estimate is
\begin{equation}
\label{eq:ch2:map_general}
\mathcal{X}^{\star} = \arg\max_{\mathcal{X}} \; p(\mathcal{X}) \prod_{k} p(\mathbf{z}_k \mid \mathcal{X}_k)
        = \arg\min_{\mathcal{X}} \; - \log p(\mathcal{X}) + \tfrac{1}{2} \sum_{k} \|\mathbf{r}_k(\mathcal{X}_k,\mathbf{z}_k)\|^2_{\bm{\Sigma}_k}.
\end{equation}
With Gaussian factors and a flat or Gaussian prior, MAP inference reduces to a nonlinear least-squares problem.

\subsection{Factor Graph Optimization}
\label{sec:ch2_factor_graph}

Factor graph optimization \cite{dellaert2017factor} provides the unifying estimation framework for this thesis.
A factor graph is a bipartite graph $\mathcal{G} = (\mathcal{X}, \mathcal{F}, \mathcal{E})$, where variable nodes $\mathcal{X} = \{\mathbf{x}_1, \ldots, \mathbf{x}_n\}$ represent unknown states (poses, velocities, biases, landmarks) and factor nodes $\mathcal{F} = \{f_1, \ldots, f_m\}$ encode probabilistic constraints derived from sensor measurements.
The MAP estimate solves the nonlinear least-squares problem
\begin{equation}
\label{eq:ch2:map}
\mathcal{X}^{\star} = \arg\min_{\mathcal{X}} \sum_{k} \left\| \mathbf{r}_k(\mathcal{X}_k) \right\|^2_{\bm{\Sigma}_k},
\end{equation}
which is solved iteratively with Gauss-Newton or Levenberg-Marquardt updates of the form
\begin{equation}
\label{eq:ch2:gn_update}
\delta \mathcal{X}^{\star} = -\left(\mathbf{J}^{\top} \mathbf{W} \mathbf{J}\right)^{-1} \mathbf{J}^{\top} \mathbf{W} \mathbf{r},
\end{equation}
where $\mathbf{J}$ stacks the residual Jacobians and $\mathbf{W}$ is the block-diagonal information matrix built from $\bm{\Sigma}_k^{-1}$.
The marginal posterior covariance at the solution follows the Laplace approximation
\begin{equation}
\label{eq:ch2:laplace}
p(\mathcal{X} \mid \mathcal{Z}) \approx \mathcal{N}(\mathcal{X}^{\star}, \mathbf{P}^{\star}), \qquad \mathbf{P}^{\star} = \left(\mathbf{J}^{\top} \mathbf{W} \mathbf{J}\right)^{-1}.
\end{equation}
The covariance $\bm{\Sigma}_k$ determines how much the optimizer trusts each measurement.
Miscalibrated covariances lead to over- or under-weighting of factors, which is the central problem this thesis addresses.

\section{SLAM and Visual-Inertial Estimation}
\label{sec:ch2_slam}

Simultaneous Localization and Mapping (SLAM) jointly estimates the trajectory of a moving sensor and a representation of its environment from sensor data, formulated as the MAP problem of \sref{sec:ch2_map} where the optimizer recovers robot poses (and, when needed, velocities, IMU biases, and a map) from odometry, observation, and loop-closure residuals weighted by their covariances.
The local-only restriction in which the map is marginalized and only consecutive poses are tracked yields odometry; adding map variables and loop closures yields full SLAM.
Visual-inertial estimation specializes this formulation by augmenting each keyframe state with velocity and IMU biases and combining visual reprojection residuals (\sref{sec:ch2_vo}) with IMU preintegration residuals (\sref{sec:ch2_imu_preint}); the visual and inertial covariances jointly determine how the optimizer balances the two modalities, and replacing them with learned, metrics-aware quantities is the core technical move of Part~III of this thesis.

\section{Stereo Visual Odometry and Camera Projection}
\label{sec:ch2_vo}

A pinhole camera with intrinsic matrix $\mathbf{K}$ projects a 3D point ${}^{c}\mathbf{p} = (X,Y,Z)^{\top}$ into the image plane through
\begin{equation}
\label{eq:ch2:projection}
\pi({}^{c}\mathbf{p}) = \mathbf{K} \begin{bmatrix} X/Z \\ Y/Z \\ 1 \end{bmatrix},
\end{equation}
and a calibrated stereo pair triangulates a 3D point ${}^{c}\mathbf{p}$ from a horizontal disparity $d$ between the left and right images using the baseline $b$ and focal length $f$, with $Z = fb/d$.

Stereo visual odometry (VO) estimates frame-to-frame camera motion from stereo image sequences.
A VO front-end establishes feature correspondences and triangulates 3D points $\mathbf{p}_i \in \mathbb{R}^3$, and a VO back-end estimates the relative pose $\mathbf{T}^{c_{k-1}}_{c_k}$ by minimizing
\begin{equation}
\label{eq:ch2:vo}
\mathbf{T}^{c_{k-1}}_{c_k}{}^{\star} = \arg\min_{\mathbf{T}} \sum_{i} \left\| \mathbf{p}^{c}_{k-1,i} - \mathbf{T} \mathbf{p}^{c}_{k,i} \right\|^2_{\bm{\Sigma}_i},
\end{equation}
where $\bm{\Sigma}_i$ is the covariance of the $i$-th correspondence and determines its influence on pose estimation.
In practice, VO robustness depends critically on this weighting: ambiguous or mismatched correspondences should be down-weighted, while geometrically consistent correspondences should dominate.
Traditional VO uses fixed diagonal covariance or disparity-only heuristics, which become unreliable under illumination change, low texture, motion blur, and dynamic objects.
This limitation motivates learning metrics-aware covariance for correspondence selection and residual weighting in Part~I (Chapter~\ref{chapter:macvo}).

\section{IMU Integration and Uncertainty Propagation}
\label{sec:ch2_imu_preint}

\subsection{IMU Kinematic Model}
\label{sec:ch2_imu_kinematics}

An inertial measurement unit (IMU) reports angular velocity ${}^{b}\bm{\omega}$ and specific force ${}^{b}\mathbf{a}$ at high rate.
The continuous-time IMU model in the body frame is
\begin{equation}
\label{eq:ch2:imu_model}
\begin{aligned}
{}^{b}\tilde{\bm{\omega}}_t &= {}^{b}\bm{\omega}_t + \mathbf{b}^g_t + \bm{\eta}^g_t,\\
{}^{b}\tilde{\mathbf{a}}_t &= \mathbf{R}_t^{\top}\!\left({}^{w}\mathbf{a}_t - {}^{w}\mathbf{g}\right) + \mathbf{b}^a_t + \bm{\eta}^a_t,
\end{aligned}
\end{equation}
with gyroscope and accelerometer biases $\mathbf{b}^g_t,\mathbf{b}^a_t$ that follow random walks and zero-mean Gaussian noises $\bm{\eta}^g_t,\bm{\eta}^a_t$ with covariances $\bm{\Sigma}^g$ and $\bm{\Sigma}^a$.
Continuous-time integration of the kinematics yields the discrete update
\begin{equation}
\label{eq:ch2:imu_kinematics}
\begin{aligned}
\mathbf{R}_{k+1} &= \mathbf{R}_k \, \mathrm{Exp}\!\left(({}^{b}\tilde{\bm{\omega}}_k - \mathbf{b}^g_k)\Delta t\right),\\
{}^{w}\mathbf{v}_{k+1} &= {}^{w}\mathbf{v}_k + {}^{w}\mathbf{g}\,\Delta t + \mathbf{R}_k({}^{b}\tilde{\mathbf{a}}_k - \mathbf{b}^a_k)\Delta t,\\
{}^{w}\mathbf{p}_{k+1} &= {}^{w}\mathbf{p}_k + {}^{w}\mathbf{v}_k\Delta t + \tfrac{1}{2}{}^{w}\mathbf{g}\,\Delta t^2 + \tfrac{1}{2}\mathbf{R}_k({}^{b}\tilde{\mathbf{a}}_k - \mathbf{b}^a_k)\Delta t^2.
\end{aligned}
\end{equation}

\subsection{Preintegration on the Manifold}
\label{sec:ch2_preint}

IMU preintegration \cite{forster2015imu} summarizes the high-rate IMU stream between two keyframes $i$ and $j$ as relative motion increments that do not require re-integration when the linearization point of the bias changes.
The preintegrated rotation, velocity, and position increments are
\begin{equation}
\label{eq:ch2:preint}
\begin{aligned}
\bm{\gamma}_{ij} &= \prod_{k=i}^{j-1} \mathrm{Exp}\!\left(({}^{b}\tilde{\bm{\omega}}_k - \mathbf{b}^g_i)\Delta t\right), \\
\bm{\beta}_{ij} &= \sum_{k=i}^{j-1} \bm{\gamma}_{ik}\,({}^{b}\tilde{\mathbf{a}}_k - \mathbf{b}^a_i)\,\Delta t, \\
\bm{\alpha}_{ij} &= \sum_{k=i}^{j-1} \left( \bm{\beta}_{ik}\,\Delta t + \tfrac{1}{2}\bm{\gamma}_{ik}\,({}^{b}\tilde{\mathbf{a}}_k - \mathbf{b}^a_i)\,\Delta t^2 \right),
\end{aligned}
\end{equation}
where $\bm{\gamma}_{ij}$, $\bm{\beta}_{ij}$, and $\bm{\alpha}_{ij}$ denote the rotation, velocity, and position preintegration terms, respectively.

\subsection{Covariance Propagation}
\label{sec:ch2_imu_cov}

The associated covariance $\bm{\Sigma}^{\mathrm{pre}}_{ij}$ over the preintegration error vector $\delta\mathbf{z}_{ij} = [\delta\bm{\phi}_{ij}^{\top},\delta\bm{\beta}_{ij}^{\top},\delta\bm{\alpha}_{ij}^{\top}]^{\top}$ is propagated through the linearized error dynamics
\begin{equation}
\label{eq:ch2:imu_cov}
\bm{\Sigma}^{\mathrm{pre}}_{i,k+1} = \mathbf{A}_k\, \bm{\Sigma}^{\mathrm{pre}}_{ik}\, \mathbf{A}_k^{\top} + \mathbf{B}^g_k\, \bm{\Sigma}^g\, (\mathbf{B}^g_k)^{\top} + \mathbf{B}^a_k\, \bm{\Sigma}^a\, (\mathbf{B}^a_k)^{\top},
\end{equation}
where $\mathbf{A}_k$ is the state-transition matrix and $\mathbf{B}^g_k$, $\mathbf{B}^a_k$ inject gyroscope and accelerometer noise into the preintegration error.
Traditional methods set $\bm{\Sigma}^g$ and $\bm{\Sigma}^a$ to fixed values from datasheet specifications or Allan deviation analysis.
Such fixed parameters fail to capture motion-dependent and platform-dependent noise characteristics, which motivates learning the IMU sensor noise and the propagation parameters in Part~II (Chapter~\ref{chapter:airimu}).

\section{Related Work}
\label{sec:ch2_related_works}

\textbf{Visual uncertainty for odometry.}
Classical visual odometry and SLAM methods rely on geometric residuals with fixed or heuristic covariance weighting \cite{mur2015orb, engel2014lsd, qin2018vins}.
Learning-based methods improve matching and depth quality \cite{sarlin2020superglue, teed2021droid, wang2023dust3r}, but most uncertainty outputs are relative confidence scores without physically meaningful scale.
The remaining gap is metrics-aware uncertainty that can be used directly in estimation and remains reliable under domain shift.
This gap motivates Part~I (Chapter~\ref{chapter:macvo}).

\textbf{Inertial uncertainty and representation.}
Traditional inertial odometry uses calibrated noise parameters and analytic propagation \cite{forster2015imu, tedaldi2014robust}, while learning-based IO improves motion estimation and bias handling \cite{chen2018ionet, herath2020ronin, liu2020tlio}.
However, prior work either assumes fixed uncertainty parameters or relies on representations that do not transfer well to high-dynamic platforms.
The key gap is jointly learning inertial uncertainty and observability-preserving representation for robust transfer across IMU grades and platforms.
This gap motivates Part~II (Chapters~\ref{chapter:airimu} and~\ref{chapter:airio}).

\textbf{Visual-inertial fusion and adaptation.}
Modern VIO and VI initialization systems are accurate but require hand-tuned weighting and calibration choices \cite{campos2021orb, qin2018vins, geneva2020openvins}.
Existing online calibration and fusion frameworks improve robustness but still depend on careful parameter adjustment \cite{huang2018online, zhao2021super}.
At the same time, learned inertial odometry models are typically trained with supervised pose or velocity labels and degrade when deployed on a new robot embodiment, mounting, or motion regime~\cite{liu2020tlio, cioffi2023learned}.
The gap is a shared, metrics-aware uncertainty representation that supports tuning-free initialization and calibration, and cross-platform adaptation of learned models under embodiment shift.
This gap motivates Part~III (Chapters~\ref{chapter:macio} and \ref{chapter:macinit}).
